% !TeX root = ../papers/nc-p3.tex
\section{判別式の他方の符号と楕円曲線上の中心}
\label{sec:klein-second-pairings}
前節の記号を用い、$D_0=ad-bc$と書く。薄い成分は全て非零で、
$z\ne\pm1$、交差ブロックと四つの組内ブロックは可逆とする。

\begin{theorem}\label{thm:klein-other-pairings}
$\delta\ne0$または$pq\ne0$なら、全ての補完は楕円曲線族の根の組合せまたは
以下の四表示で尽くされる。後者は$x=y=D_0$であり、
\[
\begin{array}{c|c}
 \lambda & A_1 \\ \hline
 1&A_0\\
 -1&-A_0\\
 1/z&\left(\begin{smallmatrix}d&c/z\\zb&a\end{smallmatrix}\right)\\
 -1/z&-\left(\begin{smallmatrix}d&c/z\\zb&a\end{smallmatrix}\right)
\end{array}
\]
となる。四表示のAS性は全て、両対称窓の既存の有限判定で決められる。
ここでは$\delta\ne0$で二次式の最高次係数が消える場合も含め、
楕円曲線側の組合せとは二つの射影化した根のうち$s,-s$に対応するものをいう。
\end{theorem}
\begin{proof}
$\delta=0$, $pq\ne0$なら各々の根は一つしかない。以下$\delta\ne0$とする。
二つの二次式を各々二変数に斉次化すると、根はそれぞれ相異なる二点である。
有限な根での符号は$2p x+m$または$2q y+k$で決まる。たとえば$p=0$なら
有限根の符号は$m$、無限遠の根の符号は$-m$である。
この対応は判別式を可逆にした全体で正則である。実際、無限遠座標$u=1/x$では
符号を$-m-2nu$と書け、共通部分で一致する。

残る四式は、分母を払うと$(x,y)$で双一次になる。たとえば対角比較の二本は
\[
\begin{split}
 (b+c-f/\lambda-\lambda g)y
   -(e+h-b-c)(x-D_0)&=0,\\
 (z^2b+c-\lambda z^2 f-g/\lambda)y
   -\{z(e+h)-z^2 b-c\}(x-D_0)&=0 .
\end{split}
\]
非対角比較の二本は、第一項がそれぞれ
$(\lambda ze+h/\lambda-za-d)x$, $(e/\lambda+\lambda zd-a-zd)x$で、
そこから$(za+d-zf-g)(y-D_0)$, $(a+zd-zf-g)(y-D_0)$を引く。
各式の双斉次化は$s,-s$の組の上で恒等的に零である。これは前節の計算と、
判別式が単元である根の有限エタール被覆で、一般点での等式が保たれることによる。
従って他方の符号の有限対$(x,y)$も解なら、各式は$x$固定で異なる二つの
射影$y$に消え、$y$固定でも異なる二つの射影$x$に消える。

まず$x$を固定した対角二式の定数係数を取る。
$x\ne D_0$なら$2e_0=b+c$, $2ze_0=z^2b+c$から$b=c=0$である。
逆に$y\ne D_0$なら非対角二式の定数係数から
$za+d=a+zd=zf_0+g_0$、すなわち$a=d=0$である。
従って$b,c$のいずれかが非零なら$x=D_0$であり、
\[
 \begin{pmatrix}1&1\\z^2&1\end{pmatrix}\binom fg
 =\begin{pmatrix}1/\lambda&\lambda\\\lambda z^2&1/\lambda\end{pmatrix}\binom bc .
\]
これは$F_1(C)$の対角条件である。さらに上の対角二式が$y$について恒等的に
零なので、$(f,g)$と$(b,c)$を交換した式も成立する。
左の係数行列の逆と右の係数行列の積の固有値は
$(1/\lambda\pm\lambda z)/(1\pm z)$である。その一つの平方が1でなければならず、
$\lambda\in\{1,-1,1/z,-1/z\}$を得る。
$b=c=0$なら$(a,d)\ne(0,0)$で、$y=D_0$を使った非対角二式が同じ結論を与える。
今度の二行列は対角$z$、非対角1の行列と、対角$\lambda z$、非対角$1/\lambda$の行列である。

これら四倍率で表の$(e,f,g,h)$を入れると$x=y=D_0$で、元の四式を満たす。
この根で$2pD_0+m=2qD_0+k$となる。これは表示した一次式への代入でも、
両辺の平方が同じことと各倍率での非零な一例からも従う。
同じ符号のもう一つの組が有限であれば、その両座標はともに$D_0$と異なるため、
前の定数係数の結論が$A_0=0$を強いてしまう。これで全被覆を得る。

$\lambda=1,A_1=A_0$は両対称である。$\lambda=1/z$の表示でも
$r=(1,1,z,1)$, $t=(1,1/z,1,1)$として
\[
 r_i w_{ijkl}=r_{\bar k}w_{\bar k,j,\bar i,l},\qquad
 t_j w_{ijkl}=t_l w_{ilkj}.
\]
従って偶相手の文字を平行移動した後に対角倍率で両対称となる。
負の二倍率は四元群の二コサイクル
$\mu(i,j)=(-1)^{i_{\rm hi} j_{\rm lo}}$による捻りで上記二例に戻せる。
実際、$H=\operatorname{diag}(1,-1)$として
$(A_0,A_1,\lambda)\mapsto(HA_0H,-HA_1H,-\lambda)$で、
厚い色の$M_h(i,k)$には$(-1)^{j_{\rm lo}+k_{\rm lo}}$が掛かり薄い色は変わらない。
この符号と四文字の$\prod_{r<s}\mu(i_r,i_s)$の差は、総次数零の各成分で
$j_{\rm lo}+k_{\rm lo}+j_{\rm hi}j_{\rm lo}+l_{\rm hi}l_{\rm lo}$の符号であり、
一文字の基底変更で吸収できる。
コサイクル恒等式により捻った乗法は結合的で、任意の斉次テンソル積に
累積文字次数による同じ符号を付ければ左右のKoszul複体の微分が対応する。
双対の完全対合も次数対ごとの非零倍率しか変わらない。よって捻りとその逆は
標準Koszul左右AS性を保つ。\path{computations/verify_klein_exceptional_pairings.py}で
上の二交換と符号を係数ごとに検算する。
\end{proof}

\begin{proposition}\label{prop:klein-elliptic-center}
$(a,b,c,d,z)=(5,2,3,7,-2)$を固定した楕円曲線の関数体上で、
四次中心の文字次数別次元は$(3,2,0,0)$である。
$\lambda=1,s=1$を含む一つの開集合の全閉点が中心四変数上で
各始点自由階数256の標準AS窓となる。
\end{proposition}
\begin{proof}
前節の四次多項式を7569で割ったものを$J(\lambda)$と書き、
$\mathbb Q(\lambda)[r]/(r^2-J(\lambda))$で$s=r/\lambda$とする。
\path{computations/verify_klein_elliptic_center.py}は各係数を
$(P(\lambda)+r Q(\lambda))/D(\lambda)$の形で四つの実際の窓を順に掃き出す。
次数5までの次元は標準値、文字次数0の36列の中心方程式の核次元は3、
次数1の40列では2となる。$\lambda=1,s=1$は
定理\ref{thm:klein-quartic-center-slice}の認証した有理点で、残る二文字次数の核は零、
かつ四中心元の商は次数13で消える。定階数の核を近傍へ延ばす前と同じ論法を適用する。
\end{proof}

\begin{corollary}[楕円曲線族の一般の印なし同型]\label{cor:klein-elliptic-isom}
四配置の射影安定群が指定の四元群に等しく、規格化と単純な根の条件を満たす
楕円曲線族の開集合で、印なし添字保持同型は高々32通りの文字置換と係数比較で決まる。
\end{corollary}
\begin{proof}
辺がグラフで、各配置の点ごとの射影固定群は自明なので、四位置の群自己同型は同じである。
文字の置換は共通の$\phi\in\operatorname{Aut}(K)$に四つの平行移動量
$u,v,u',v'$を加えたものとなる。総次数と両方の厚い次数の組から
$u+v+u'+v'=0$, $\phi(1)=1$, $u+v,u'+v\in\{0,1\}$を得る。
従って$2\cdot4\cdot2\cdot2=32$通りである。各置換後に薄い係数を規格化し、
$A_0$の一つの非零成分で両方の$A_0,A_1$を割る。$z,\lambda$とこの二行列の一致が
必要十分である。実際、残る薄い次数の逆倍率の四乗根の選び方を変えても、
各組内での行・列の相対倍率は同一の符号に限られ、両行列は同じスカラー倍しか変化しない。
逆にこれらの値は全テンソルを再構成する。適切な基底成分が零なら別の非零成分で割ればよい。
\end{proof}
